\labsection{1}{Feature engineering and feature selection}{sec:lab01-features}

\begin{labproject}{Check the tools, understand the data, build features, and choose among them}
\textbf{Goal.} Check the environment, engineer features, compare scaling, and select features without leakage.

\medskip\noindent\textbf{Setup.}\par
\begin{lstlisting}[style=mlpython]
from pathlib import Path

import matplotlib
import matplotlib.pyplot as plt
import numpy as np
import pandas as pd
import sklearn
from sklearn.compose import ColumnTransformer
from sklearn.feature_selection import SelectKBest, f_regression
from sklearn.impute import SimpleImputer
from sklearn.linear_model import Lasso, LogisticRegression, Ridge
from sklearn.metrics import mean_absolute_error, pairwise_distances, roc_auc_score
from sklearn.model_selection import train_test_split
from sklearn.pipeline import Pipeline
from sklearn.preprocessing import OneHotEncoder, StandardScaler

# Notebooks live in all_experiments/, so the course folder is one level up.
PROJECT_ROOT = Path.cwd().parent
DATA = PROJECT_ROOT / "all_datasets"
\end{lstlisting}

\medskip\noindent\textbf{Part A --- Environment check and a first look at the data.}\par

Record package versions for reproducibility.

\begin{lstlisting}[style=mlpython]
print("NumPy:", np.__version__)
print("pandas:", pd.__version__)
print("Matplotlib:", matplotlib.__version__)
print("scikit-learn:", sklearn.__version__)
\end{lstlisting}

Titanic has one row per passenger; \texttt{Survived}=1 denotes survival. Inspect shape, types, missingness, and class balance.

\begin{lstlisting}[style=mlpython]
df = pd.read_csv(DATA / "titanic_dataset.csv")

print("shape:", df.shape)
print(df.head())
print("\nData types:\n", df.dtypes)
print("\nSurvival rate:\n", df["Survived"].value_counts(normalize=True))
\end{lstlisting}

Plot missing-value percentages by column.

\begin{lstlisting}[style=mlpython]
missing = df.isna().mean().sort_values(ascending=False)
missing = missing[missing > 0]

plt.figure(figsize=(5, 3))
plt.barh(missing.index, missing.values * 100)
plt.xlabel("Share of rows missing (%)")
plt.title("Titanic: where the gaps are")
plt.tight_layout()
plt.show()
\end{lstlisting}

\medskip\noindent\textbf{Part B --- Feature engineering.}\par

Derive family size, isolation, title, and log fare from prediction-time information. Each feature must exist before the outcome.

\begin{lstlisting}[style=mlpython]
df["FamilySize"] = df["SibSp"] + df["Parch"] + 1
df["IsAlone"] = (df["FamilySize"] == 1).astype(int)
df["Title"] = df["Name"].str.extract(r",\s*([^\.]+)\.")[0].str.strip()
df["Title"] = df["Title"].where(df["Title"].isin(["Mr", "Mrs", "Miss", "Master"]), "Other")
df["LogFare"] = np.log1p(df["Fare"])

print(df[["Name", "Title", "FamilySize", "IsAlone", "Fare", "LogFare"]].head())
print("\nTitle counts:\n", df["Title"].value_counts())
\end{lstlisting}

Compare survival rates across engineered features.

\begin{lstlisting}[style=mlpython]
fig, axes = plt.subplots(1, 2, figsize=(9.5, 3.4))
df.groupby("Title")["Survived"].mean().sort_values().plot.barh(ax=axes[0])
axes[0].set_xlabel("Survival rate")
axes[0].set_title("Title carries the sex-and-age story in one column")
df.groupby("FamilySize")["Survived"].mean().plot.bar(ax=axes[1], rot=0)
axes[1].set_ylabel("Survival rate")
axes[1].set_title("Small families did best; large ones did worst")
plt.tight_layout()
plt.show()
\end{lstlisting}

\textbf{Inspect.} Survival varies by title and family size; small families fare better than lone passengers or large families.

Compare original and engineered features using the same model and split. Fit imputation, scaling, and encoding inside the training pipeline.

\begin{lstlisting}[style=mlpython]
def make_preprocess(num_cols, cat_cols):
    return ColumnTransformer([
        ("num", Pipeline([
            ("impute", SimpleImputer(strategy="median")),
            ("scale", StandardScaler()),
        ]), num_cols),
        ("cat", Pipeline([
            ("impute", SimpleImputer(strategy="most_frequent")),
            ("onehot", OneHotEncoder(handle_unknown="ignore")),
        ]), cat_cols),
    ])

y = df["Survived"]
feature_sets = {
    "raw columns": (["Age", "SibSp", "Parch", "Fare"], ["Pclass", "Sex", "Embarked"]),
    "raw + engineered": (["Age", "SibSp", "Parch", "LogFare", "FamilySize", "IsAlone"],
                         ["Pclass", "Sex", "Embarked", "Title"]),
}

for name, (num_cols, cat_cols) in feature_sets.items():
    X = df[num_cols + cat_cols]
    X_train, X_valid, y_train, y_valid = train_test_split(
        X, y, test_size=0.25, random_state=42, stratify=y
    )
    model = Pipeline([
        ("preprocess", make_preprocess(num_cols, cat_cols)),
        ("model", LogisticRegression(max_iter=2000)),
    ]).fit(X_train, y_train)
    auc = roc_auc_score(y_valid, model.predict_proba(X_valid)[:, 1])
    print(f"{name:18s} validation ROC AUC = {auc:.3f}")
\end{lstlisting}

\textbf{Inspect.} Higher validation AUC measures the gain from features, not a different estimator.

\medskip\noindent\textbf{Part C --- Scaling changes the geometry.}\par

Breast-cancer measurements have unequal scales: area reaches thousands while smoothness is below one. Compare their effect on distances.

\begin{lstlisting}[style=mlpython]
cancer = pd.read_csv(DATA / "breast_cancer_dataset.csv")
feature_cols = [c for c in cancer.columns if c not in {"id", "diagnosis", "Unnamed: 32"}]
X_cancer = cancer[feature_cols]

three = ["area_mean", "texture_mean", "smoothness_mean"]
print(X_cancer[three].describe().loc[["mean", "std"]].round(3))
\end{lstlisting}

Compare raw and standardized distances among the first five tumors; darker heatmap cells indicate closer pairs.

\begin{lstlisting}[style=mlpython]
raw = pairwise_distances(X_cancer.iloc[:5], metric="euclidean")
X_scaled = StandardScaler().fit_transform(X_cancer)
scaled = pairwise_distances(X_scaled[:5], metric="euclidean")

fig, axes = plt.subplots(1, 2, figsize=(9, 3.6))
for ax, matrix, title in [(axes[0], raw, "Raw units"), (axes[1], scaled, "Standardized")]:
    im = ax.imshow(matrix, cmap="viridis")
    ax.set_title(title)
    ax.set_xlabel("tumour")
    ax.set_ylabel("tumour")
    for i in range(5):
        for j in range(5):
            ax.text(j, i, f"{matrix[i, j]:.0f}", ha="center", va="center", color="white", fontsize=8)
    fig.colorbar(im, ax=ax, shrink=0.8)
plt.suptitle("Same five tumours, two geometries")
plt.tight_layout()
plt.show()
\end{lstlisting}

\textbf{Inspect.} The closest pair changes from tumors 3/5 to 2/5 after scaling. Units can change neighbors.

\medskip\noindent\textbf{Part D --- Feature selection without leakage.}\par

Predict \texttt{G3} for 649 students at enrolment. Exclude later grades \texttt{G1} and \texttt{G2}, leaving 30 candidate columns.

\begin{lstlisting}[style=mlpython]
students = pd.read_csv(DATA / "student_performance_dataset.csv")
print("shape:", students.shape, " missing:", students.isna().sum().sum(),
      " duplicates:", students.duplicated().sum())

target = "G3"
leaky = ["G1", "G2"]
X_s = students.drop(columns=[target] + leaky)
y_s = students[target]
s_num = X_s.select_dtypes(include="number").columns.tolist()
s_cat = X_s.select_dtypes(exclude="number").columns.tolist()
print(len(s_num), "numeric and", len(s_cat), "text features available at enrolment")

Xs_train, Xs_valid, ys_train, ys_valid = train_test_split(X_s, y_s, test_size=0.25, random_state=42)
\end{lstlisting}

Compare ridge using all features or ten \texttt{SelectKBest} features, and lasso's sparse fit. Keep selection inside training pipelines. A fourth run deliberately adds \texttt{G2} to expose leakage.

\begin{lstlisting}[style=mlpython]
def evaluate(pipeline, X_tr, X_va, label):
    pipeline.fit(X_tr, ys_train)
    mae = mean_absolute_error(ys_valid, pipeline.predict(X_va))
    print(f"{label:32s} validation MAE = {mae:.3f} grade points")
    return mae

results = {}
results["all 30 features"] = evaluate(Pipeline([
    ("preprocess", make_preprocess(s_num, s_cat)),
    ("model", Ridge(alpha=1.0)),
]), Xs_train, Xs_valid, "all 30 features")

results["SelectKBest, k=10"] = evaluate(Pipeline([
    ("preprocess", make_preprocess(s_num, s_cat)),
    ("select", SelectKBest(f_regression, k=10)),
    ("model", Ridge(alpha=1.0)),
]), Xs_train, Xs_valid, "SelectKBest, k=10")

lasso = Pipeline([
    ("preprocess", make_preprocess(s_num, s_cat)),
    ("model", Lasso(alpha=0.1)),
])
results["lasso (selects itself)"] = evaluate(lasso, Xs_train, Xs_valid, "lasso (selects itself)")

leak_num = s_num + ["G2"]
results["all 30 + G2 (leakage)"] = evaluate(Pipeline([
    ("preprocess", make_preprocess(leak_num, s_cat)),
    ("model", Ridge(alpha=1.0)),
]), Xs_train.assign(G2=students.loc[Xs_train.index, "G2"]),
    Xs_valid.assign(G2=students.loc[Xs_valid.index, "G2"]), "all 30 + G2 (leakage)")
\end{lstlisting}

\begin{lstlisting}[style=mlpython]
coefs = pd.Series(lasso.named_steps["model"].coef_,
                  index=lasso.named_steps["preprocess"].get_feature_names_out())
kept = coefs[coefs != 0].sort_values()
print(f"lasso kept {len(kept)} of {len(coefs)} encoded columns")

fig, axes = plt.subplots(1, 2, figsize=(11, 3.8))
axes[0].barh(list(results), list(results.values()))
axes[0].set_xlabel("Validation MAE (grade points)")
axes[0].set_title("Fewer features, same error; a leaked feature, half the error")
axes[1].barh(kept.index, kept.values)
axes[1].axvline(0, color="grey", linewidth=0.8)
axes[1].set_title("What the lasso kept")
plt.tight_layout()
plt.show()
\end{lstlisting}

\textbf{Inspect.} Selected features give similar or better error with fewer columns. Adding \texttt{G2} cuts error sharply but invalidates an enrolment-time forecast. Inspect lasso's retained school, failure, education, alcohol, and study-time features without treating them as causal.

\begin{tryit}
Add cabin initial or fare per family member; compare Part~B validation AUC. Rerun Part~D with \texttt{k=5} and \texttt{k=20}; compare errors.
\end{tryit}
\end{labproject}
